\documentclass[10pt,a4paper]{amsart}
\usepackage[margin=19mm]{geometry}
\usepackage{amsmath,amssymb,mathtools}
\usepackage[scr=rsfs,cal=euler]{mathalfa}
\usepackage{xfrac}
\usepackage[unicode,hidelinks,bookmarks=false]{hyperref}
\newcommand{\ie}{i.e.\ }
\newcommand{\cf}{cf.\ }
\newcommand{\resp}{respectively\ }
\newcommand{\Subsection}{\subsection}
\providecommand{\nobreakdash}{\nobreak\hbox{-}}
\setlength{\emergencystretch}{2em}
\multlinegap0pt
\usepackage{bm}
\usepackage{bbm}
\usepackage{dsfont}
\usepackage{xspace}
\usepackage{cancel}
\usepackage{mathtools}
\usepackage{cleveref}
\usepackage{amsthm}
\makeatletter
\@ifundefined{theorem}{\newtheorem{theorem}{Theorem}}{}
\makeatother
\title[On the maximum number of edges of outer k-planar graphs]{On the maximum number of edges of outer k-planar graphs}
\author{Maximilian Pfister}
\date{Source: arXiv:2505.24490v1 (CC BY 4.0)}
\begin{document}
\maketitle

\noindent\emph{Source and licence.} On the maximum number of edges of outer k-planar graphs. Version arXiv:2505.24490v1 (CC BY 4.0), distributed under the Creative Commons Attribution 4.0 International licence, as recorded in the arXiv licence metadata for this exact version and shown on the abstract page https://arxiv.org/abs/2505.24490v1. Full rights notice: licenses/arxiv-2505.24490-CC-BY-4.0.txt.\par\smallskip

\noindent\textit{Editorial scope.} Counted unit: the Theorem whose source label is \texttt{thm:upper-new} in the article (source0.tex lines 281--284, source label \texttt{thm:upper-new}), delivered together with the article's own complete original proof of it (source0.tex lines 285--309, one \texttt{proof} environment of 25 lines). No mathematics is written, repaired or re-derived: statement and proof are the article's own text line for line. Required context retained uncounted: none. Every definition, display and label of those lines is retained unchanged. Omitted from this package and not redistributed: 1--280, 310--534. Nothing inside a retained excerpt is omitted or altered: every retained body is delivered exactly as the article's lines are written, so no omission receipt of any kind accompanies this package. Citations: the retained excerpts carry no citation command at all, so no bibliography entry of the article is transcribed and no reference is redistributed. Reference adaptation: none was needed and none was made; every reference inside the delivered unit resolves inside it, so no unresolved reference or citation is produced. Printed slips kept literal: no printed word, symbol, hypothesis or number of the counted statement or of its proof was altered. Layout and class: the article's own class is \texttt{socg} with packages including inputenc, amsmath, amssymb, todonotes, enumerate, xspace, xcolor, soul, booktabs, tabularx, colortbl, adjustbox, comment, graphicx; the wrapper is the lane's pinned \texttt{amsart} editorial wrapper at 10pt on A4, which restates the theorem-like environments the retained text needs and adds the article's own macro definitions that the retained text needs (none), with extra wrapper packages (bm, bbm, dsfont, xspace, cancel, mathtools, cleveref). The delivered numbering is the unit's own. Assets: no retained span contains a figure environment, a graphics-inclusion command, a PDF-page inclusion, a table or a TikZ picture; no image file is copied into this package, the package declares no asset dependency and no image file is redistributed with it. Licence: version arXiv:2505.24490v1 of this article is distributed under the Creative Commons Attribution 4.0 International licence, as recorded in the arXiv licence metadata for this exact version and shown on the abstract page https://arxiv.org/abs/2505.24490v1. A full rights notice is delivered at licenses/arxiv-2505.24490-CC-BY-4.0.txt. Characters: every retained excerpt is pure ASCII, so the delivered engine, XeLaTeX with its default Latin Modern text font, reports no missing character.
\begin{theorem} \label{thm:upper-new}
An outer $k$-planar graph has at most $(\sqrt{2}+\varepsilon)\sqrt{k}n$ diagonals with 
$\lim_{k\to \infty} \varepsilon= 0$.
\end{theorem}

\begin{proof}
Let $x = (\sqrt{2}+\varepsilon)$.
Let $G$ be a vertex minimum counterexample to \cref{thm:upper-new} and let $v_1,\dots,v_n$ be the vertices of $G$, where the order is defined by the outer face. Fix a parameter $1 \leq l_0 \leq \frac{n}{2}$ and choose a shortest diagonal of $G$ of length at least $l_0$. Denote the diagonal by $D$, assume its length is $l \geq l_0$ and w.l.o.g. assume that $D = v_1v_{l+3}$. Let $G_1$ and $G_2$ denote the two outer $k$-planar graphs obtained from $G$ by splitting along $D$ and removing all diagonals which intersect $D$. W.l.o.g. assume that $G_1$ contains $l+2$ vertices on its boundary, which implies that $G_2$ contains $n-l+2$ vertices.
Recall that by our choice of $D$, any diagonal contained in $G_1$ has length at most $l_0$.
Note that the length of a diagonal is meassured w.r.t. $G$, that is, without interpreting $D$ as a boundary edge. Thus, the number of diagonals of length two contained in $G_1$ is exactly $l$, since the two diagonals $v_{l+2}v_1$ and $v_{l+3}v_2$ do not have length two w.r.t. $G$. 
Likewise, the number of diagonals of length three is $l-1$, as the three diagonals $\{v_{l+1}v_1,v_{l+2}v_2,v_{l+3}v_3\}$ do not have length three w.r.t. $G$. In general, there exist at most $(l+2-p)$ many diagonals of length $p$ in $G_1$.
Hence, the number of diagonals in $G_1$ is at most
\[(l+2-2)+(l+2-3)+\dots(l+2-l_0) = (l_0-1)\frac{(2l-l_0+2)}{2}\]
Since $G$ is chosen as a minimum counterexample, it follows that the number of diagonals that are contained in $G_2$ is at most $x(n-l+2)$.
This implies that $D$ is crossed by at least
\[\underbrace{xn+1}_{\text{G counterexample}} -\underbrace{1}_D - \underbrace{x(n-l+2)}_{G_2} - \underbrace{(l_0-1)\frac{(2l-l_0+2)}{2}}_{\text{Diagonals in }G_1} = xl-2x - (l_0-1)\frac{(2l-l_0+2)}{2} \] edges.
In order to obtain a contradiction, set
\[xl-2x - (l_0-1)\frac{(2l-l_0+2)}{2} \stackrel{!}{>} k\]
\[\Leftrightarrow xl-2x-l_0l+0.5l_0^2-l_0+l-0.5l_0+1 \stackrel{!}{>}k\]

Choose $l_0 = \sqrt{2}\sqrt{k}$. Recall that $x = (\sqrt{2}+\varepsilon)\sqrt{k}$ and that $l \geq l_0$ holds. 
This implies that $xl-l_0l = \varepsilon\sqrt{k}l$. Since $0.5l_0^2 = k$, we have

\[k + 1 + \varepsilon\sqrt{k}l -2(\sqrt{2}+\varepsilon)\sqrt{k} -0.5\sqrt{2}\sqrt{k} \stackrel{!}{>}k \]
Thus, we have to assert
\[\varepsilon\sqrt{k}l + 1 \stackrel{!}{>} 2(\sqrt{2}+\varepsilon)\sqrt{k} +0.5\sqrt{2}\sqrt{k}\]
Since $l \geq l_0 = \sqrt{2}\sqrt{k}$, the LHS is at least $\varepsilon\sqrt{k}(\sqrt{2}\sqrt{k})+1$ and thus
\[(\sqrt{2}k-2\sqrt{k})\varepsilon \stackrel{!}{>}\frac{5\sqrt{2}}{2}\sqrt{k}-1 \]
Since the LHS side grows linearly in $k$, while the RHS has a sublinear growth, there exists, for any $\varepsilon > 0$, a sufficiently large $k$ for which the LHS is larger than the RHS as desired.
\end{proof}

\end{document}
